\begin{abstract}
The proliferation of Unmanned Aerial Vehicles (UAVs) in cooperative Flying Ad-hoc Networks (FANETs) has underscored the critical need for secure, decentralized, and efficient communication systems. Traditional FANETs remain vulnerable to cyber-physical attacks and are often hindered by the single-point-of-failure risks associated with centralized architectures. This paper presents a comprehensive analysis of blockchain technology as a foundational solution to these challenges. We introduce the PrimeFusion FANET framework, a novel, multi-layered architecture designed specifically for the resource-constrained environment of UAVs. The PrimeFusion framework integrates three key innovations: an efficient CBOR-based data compression layer to minimize bandwidth usage, a lightweight Proof-of-Cooperation (PoC) consensus mechanism that incentivizes collaborative behavior while reducing computational overhead, and a robust, multi-faceted security system. This work provides a critical evaluation of the proposed framework in comparison to both traditional FANET protocols and generic, computationally-intensive blockchain solutions. Our analysis demonstrates that an adapted blockchain implementation can significantly enhance data integrity, trust, and resilience in UAV networks. However, a balanced discussion also addresses the inherent limitations and trade-offs of this approach, including residual latency, energy consumption, and scalability constraints. We conclude that while not a universal solution, tailored blockchain frameworks like PrimeFusion represent a transformative step towards enabling secure, autonomous, and truly decentralized cooperative UAV operations.
\end{abstract}
